\section{Orchestration and pairing}
\label{system-architecture-and-implementation:orchestration-and-pairing}

% LEAD: how the pairing protocol of
%   Section~\ref{research-design-and-methodology:overview-and-experimental-design:replication-and-randomization}
%   is realized as a serial-outer / parallel-inner orchestrator. Do not re-argue why pairing matters.

\subsection{Dispatch loop and batch graph}
\label{system-architecture-and-implementation:orchestration-and-pairing:dispatch-loop-and-batch-graph}

The orchestrator in \texttt{main.py} drives the suite through a single serial loop, so that only one batch of experiments is ever in flight. It walks the experiment list in order, skips any entry already recorded in \texttt{completed\_experiments}, and treats the next remaining entry as the seed of a batch. The members of that batch are the experiments named in the seed's \texttt{related\_script\_ids} list, which every peer declares reciprocally, so the union of these reciprocal declarations forms the batch graph. The seed and its peers launch together across a thread pool capped at ten workers (\texttt{ThreadPoolExecutor} with \texttt{max\_workers=10}), one Azure Container Instance per member; the loop blocks until the whole batch has finished, marks every member completed, and only then advances. Because each member is appended to \texttt{completed\_experiments} once its batch returns, the reciprocal declarations are de-duplicated and each batch executes exactly once, independent of which member the outer loop reached first. Running the peers of a batch within one wall-clock window is what makes their measurements comparable, for the reasons established in Section~\ref{research-design-and-methodology:overview-and-experimental-design:replication-and-randomization}.

Failure is contained at the level of the individual member rather than the batch. Each member runs inside \texttt{\_safe\_run}, which catches any exception raised while launching or monitoring that container, records the failed identifier, and lets the remaining peers run to completion instead of aborting the batch. After the batch returns, the orchestrator inspects the collected failures and logs whether the batch failed in part or in full, noting that surviving peers ran without their matched counterparts and that the fair-comparison assumption may not hold for that batch; the warning lets such batches be filtered out during analysis. Every member executes in its own container group provisioned with \texttt{restart-policy Never}, so a member that exits is never silently retried on the same instance and each measurement reflects a single clean execution.

\subsection{Batching constraints}
\label{system-architecture-and-implementation:orchestration-and-pairing:batching-constraints}
% [NEW] Four simultaneous constraints (CLAUDE.md / README): (a) same query type, (b) same
%   dataset_size, (c) at most one PostGIS member (shared Azure Postgres server), (d) Databricks
%   cluster vCPU sum <= 200, computed as (workers+1)x4 per Sedona member on Standard_D4s_v3 =
%   regional quota. DuckDB/Shapefile draw from a separate ACI quota. Each maps to a real shared
%   resource being protected.

\subsection{Batch listing}
\label{system-architecture-and-implementation:orchestration-and-pairing:batch-listing}
% [NEW] Reproduce the 17-batch table from README (P1,P2,K1,K2,B1,B2,A_S1,A_S2,A_M1..A_M3,
%   A_L1..A_L6) with per-batch Databricks vCPU sum and members. RESOLVE THE COUNT DISCREPANCY:
%   README says "51 experiments / 17 batches" in some places and "17 batches" with a body text
%   that elsewhere says 20; benchmarks.yml is authoritative — derive counts from it and flag any
%   mismatch with the README prose. Use booktabs; \label the table.
\begin{table}[H]
    % \centering / \begin{tabular} ... batch | type | size | databricks vCPU | members
    \caption[Benchmark batch listing.]{...}
    \label{tab:batch-listing}
\end{table}